\chapter{Технологическая часть}

\section{Реализация алгоритмов}

Для реализации алгоритмов использовался язык программирования C++, так как он поддерживает статическую типизацию и имеет возможность отключения сборщика мусора. Отладка происходила в среде разработки <<Visual Studio 2022>>.

% 17896

\lstinputlisting[firstline=137, lastline=185, style=cpp]{../code/sales.cpp}
Листинг 1 -- Реализация с полным перебором.

\lstinputlisting[linerange={5-87, 112-133}, style=cpp]{../code/sales.cpp}
Листинг 2 -- Реализация с муравьиным алгоритмом.

\section{Тестирование}

Результаты тестирования представлены в таблице 3.1.

\begin{table}[!htbp]
	\caption{Результаты тестирования.}
	\centering
	\begin{tabularx}{\textwidth}{|*{3}{>{\centering\arraybackslash}X|}}
		\hline
		\textbf{Алгоритм} & \textbf{Вход} & \textbf{Выход} \\\hline
		Полный перебор & 
		$\begin{pmatrix}
		0&2&9&10&7\\
		1&0&6&4&3\\
		9&7&0&8&5\\
		6&3&2&0&4\\
		3&8&5&6&0\\
		\end{pmatrix}$ & 0 1 3 2 4 0 \\\hline
		Муравьиный & 
		$\begin{pmatrix}
		0&2&9&10&7\\
		1&0&6&4&3\\
		9&7&0&8&5\\
		6&3&2&0&4\\
		3&8&5&6&0\\
		\end{pmatrix}$ & 4 0 1 3 2 4 \\\hline
	\end{tabularx}
\end{table}
\FloatBarrier

Начало и конец найденного муравьиным алгоритмом пути совпадают, следовательно результат верный.

% надо сравнивать разницу длин пути/что начало и конец совпадает

\section*{Вывод}
%<что делали и получили в результате>

В данном разделе мной были реализованы описанные ранее алгоритмы и проведено их успешное тестирование.

\clearpage
